\appendix

\section{Proof Details}

\subsection{Complete Responses and Pre-Instance Coordinates}
\label{app:step-001}

\begin{lemma}[Complete public response rules and the zero-query convention]
\label{lem:step-001-complete-rules}
Under Assumptions~\ref{assump:parameter-regime} and
\ref{assump:adaptive-sq-interface}, let \(\mathscr Q_A\) be the set of all
public query-bearing histories
\[
\eta_t=(q_1,v_1,\ldots,q_{t-1},v_{t-1},q_t),
\qquad 1\le t\le m,
\]
admitted by the fixed learner for some seed and some preceding public reply
sequence in \(I_\tau\), including histories not reached for a particular seed
or reply sequence.  Then
\[
\mathfrak R_A^{\mathrm{all}}=I_\tau^{\mathscr Q_A}
=\{R:\mathscr Q_A\to I_\tau\}
\]
is exactly the family of deterministic complete public-transcript response
rules and is nonempty.  Every \(R\in\mathfrak R_A^{\mathrm{all}}\) and
\(u\in\Omega_A\) determine a terminal predictor
\(g_{u,R}:\mathcal X\to\{-1,+1\}\).  If \(m=0\), then
\(\mathscr Q_A=\varnothing\) and
\(\mathfrak R_A^{\mathrm{all}}=\{R_\varnothing\}\), where
\(R_\varnothing\) is the unique empty response rule.
\end{lemma}

\begin{proof}
The set \(\mathscr Q_A\) is determined solely by the fixed public interaction
interface: its elements contain public queries and replies but no hidden-seed
coordinate.  A total function on \(\mathscr Q_A\) therefore observes the seed
only through information already present in the public transcript.  Because
its domain contains every admitted query-bearing history, including those not
reached on one seed path, such a function is complete.  Conversely, every
deterministic complete response rule assigns one value in \(I_\tau\) to every
element of \(\mathscr Q_A\), proving the equality of the two families.

Assumption~\ref{assump:parameter-regime} gives \(\tau>0\), and hence
\[
0\in[-1-\tau,1+\tau]=I_\tau.
\]
Thus the assignment \(R^{\mathrm{zero}}(\eta)=0\) for all
\(\eta\in\mathscr Q_A\) proves nonemptiness.  This rule need not be valid for
a later learning instance; it is only an all-rule witness.

For fixed \(u\) and \(R\), the learner selects its next action from \(u\) and
the preceding public transcript.  Whenever it issues a query, totality of
\(R\) supplies the next reply.  Assumption~\ref{assump:adaptive-sq-interface}
makes this recursion well defined and forces stopping after at most \(m\)
queries with a binary predictor \(g_{u,R}\).

If \(m=0\), there is no index \(t\) with \(1\le t\le m\), so
\(\mathscr Q_A=\varnothing\).  There is exactly one function from the empty
set to \(I_\tau\), namely \(R_\varnothing\).  The learner makes no query, but
its seed-dependent terminal output remains well defined.
\end{proof}

\begin{proposition}[Finite seed-averaged response space and its exact scope]
\label{prop:step-001-mean-response-space}
Under Assumptions~\ref{assump:parameter-regime},
\ref{assump:adaptive-sq-interface}, and
\ref{assump:mean-response-rank}, and Lemma~\ref{lem:step-001-complete-rules},
define
\[
F_R(x):=\int_{\Omega_A}g_{u,R}(x)\,\mu_A(du),
\quad
V_A:=\operatorname{span}_{\mathbb R}
\{F_R:R\in\mathfrak R_A^{\mathrm{all}}\},
\quad
r_A:=\dim V_A.
\]
Then every \(F_R\) is a well-defined member of
\([-1,1]^{\mathcal X}\), and
\[
r_A\le B\bigl(1+m/\tau^2\bigr)^k<\infty.
\]
This average-before-span construction gives no bound by \(r_A\) on the span
of the seed-specific functions \(g_{u,R}\).  If
\(\mathcal X=\varnothing\), then \(V_A=\{0\}\) and \(r_A=0\).  If
\(m=0\), then
\[
V_A=\operatorname{span}_{\mathbb R}\{F_{R_\varnothing}\},
\qquad r_A\in\{0,1\},
\qquad r_A\le B.
\]
\end{proposition}

\begin{proof}
Lemma~\ref{lem:step-001-complete-rules} supplies a terminal binary function
\(g_{u,R}\) for every seed and complete rule.  The measurability clause of
Assumption~\ref{assump:adaptive-sq-interface} makes the displayed seed
integral well defined.  Since \(-1\le g_{u,R}(x)\le1\), monotonicity of
integration against the probability law \(\mu_A\) gives
\[
-1\le\int_{\Omega_A}g_{u,R}(x)\,\mu_A(du)\le1.
\]
Thus \(F_R:\mathcal X\to[-1,1]\), and finite real linear combinations of
these functions form \(V_A\subseteq\mathbb R^{\mathcal X}\).

Assumption~\ref{assump:mean-response-rank} applies to exactly this space:
\[
r_A=\dim\operatorname{span}_{\mathbb R}
\left\{x\mapsto\int_{\Omega_A}g_{u,R}(x)\,\mu_A(du):
R\in\mathfrak R_A^{\mathrm{all}}\right\}
\le B\bigl(1+m/\tau^2\bigr)^k.
\]
The right-hand side is finite because \(B,k\) are finite,
\(m\in\mathbb N_0\), and \(\tau>0\).  Hence \(r_A\in\mathbb N_0\).

To make the scope of this certificate explicit, define
\[
W_A^{\mathrm{seed}}
:=\operatorname{span}_{\mathbb R}
\{g_{u,R}:u\in\Omega_A,\ R\in\mathfrak R_A^{\mathrm{all}}\}.
\]
The two operations are genuinely different:
\[
V_A=\operatorname{span}_{R}
\left\{\int g_{u,R}\,d\mu_A(u)\right\},
\qquad
W_A^{\mathrm{seed}}=\operatorname{span}_{u,R}\{g_{u,R}\}.
\]
No individual \(g_{u,R}\) is claimed to lie in \(V_A\), and no finiteness
or cardinality assertion about \(W_A^{\mathrm{seed}}\) is used below.

The distinction can be arbitrarily large even under all four assumptions.
Fix \(N\ge5\), take \(\mathcal X_N=\{1,\ldots,N\}\), \(m=0\),
\(B=k=1\), arbitrary \(\tau>0\), and \(\varepsilon=1/N<1/4\).  Let
\(\mathcal H=\{h_+\}\), where \(h_+(i)=+1\), let \(U\) be uniform on
\(\{1,\ldots,N\}\), and let the zero-query learner with seed \(j\) return
\[
g_{j,R_\varnothing}(i)=
\begin{cases}
-1,&i=j,\\
+1,&i\ne j.
\end{cases}
\]
The unique mean response is the nonzero constant
\[
F_{R_\varnothing}(i)=\frac{N-1}{N}-\frac1N=1-\frac2N,
\]
so \(r_A=1\) and the mean-response-rank certificate holds with equality.
The \(N\) terminal predictors are nevertheless linearly independent.  If
\(\sum_{j=1}^N c_jg_{j,R_\varnothing}=0\) and
\(S_c=\sum_{j=1}^Nc_j\), evaluation at \(i\) gives
\(S_c-2c_i=0\) for every \(i\).  Hence \(c_i=S_c/2\) and
\(S_c=NS_c/2\).  Since \(N\ne2\), this forces \(S_c=0\) and all
\(c_i=0\).  Thus \(\dim W_A^{\mathrm{seed}}=N\) while \(r_A=1\).

This example also satisfies Assumption~\ref{assump:universal-adversarial-guarantee}.
For every probability law \(\mathcal D\) on \(\mathcal X_N\), the only
response rule is \(R_\varnothing\), and
\[
\mathbb E_U\mathcal L_{\mathcal D,h_+}(g_{U,R_\varnothing})
=\frac1N\sum_{j=1}^N\mathcal D(\{j\})
=\frac1N=\varepsilon.
\]
Thus the gap between the two spans is compatible with the complete theorem
contract.

If \(\mathcal X=\varnothing\), there is only one real-valued function on
\(\mathcal X\), namely the zero vector of \(\mathbb R^{\mathcal X}\).
Every \(g_{u,R}\) and \(F_R\) is this function, so \(V_A=\{0\}\) and
\(r_A=0\).  If \(m=0\), Lemma~\ref{lem:step-001-complete-rules} gives the
single rule \(R_\varnothing\), so \(V_A\) has one generator and dimension
zero or one.  The primitive bound specializes exactly to
\[
r_A\le B\bigl(1+0/\tau^2\bigr)^k=B.
\]
\end{proof}

\begin{proposition}[One pre-instance basis and exact coordinate map]
\label{prop:step-001-fixed-coordinates}
Under Assumptions~\ref{assump:parameter-regime},
\ref{assump:adaptive-sq-interface}, and
\ref{assump:mean-response-rank}, and
Proposition~\ref{prop:step-001-mean-response-space}, there is a basis
\(\psi_1,\ldots,\psi_{r_A}\) of \(V_A\), with the empty basis when
\(r_A=0\), such that
\[
\varphi_A(x):=(\psi_1(x),\ldots,\psi_{r_A}(x))
\]
is fixed before every distribution, target, selected complete or valid
response rule, and realized seed.  Every \(f\in V_A\) has a unique coordinate
vector \(a(f)\in\mathbb R^{r_A}\) satisfying
\[
f(x)=\langle a(f),\varphi_A(x)\rangle
\qquad\text{for every }x\in\mathcal X.
\]
The statement includes \(r_A=0\), empty \(\mathcal X\) or \(\mathcal H\),
and \(m=0\).
\end{proposition}

\begin{proof}
Proposition~\ref{prop:step-001-mean-response-space} makes \(V_A\)
finite-dimensional with dimension \(r_A\).  Fix a linearly independent
spanning list \(\psi_1,\ldots,\psi_{r_A}\); if \(r_A=0\), the list is empty
and \(V_A=\{0\}\).

For \(r_A>0\), spanning gives coefficients
\(a(f)_1,\ldots,a(f)_{r_A}\) with
\[
f=\sum_{j=1}^{r_A}a(f)_j\psi_j.
\]
Linear independence makes the coefficients unique.  Evaluation at
\(x\in\mathcal X\) gives the zero-residual identity
\[
f(x)=\sum_{j=1}^{r_A}a(f)_j\psi_j(x)
=\left\langle a(f),(\psi_1(x),\ldots,\psi_{r_A}(x))\right\rangle
=\langle a(f),\varphi_A(x)\rangle.
\]

The dependence order is
\[
\begin{aligned}
&(\mathcal X,A,\Omega_A,\mu_A,m,\tau,
  \text{complete public interface})\\
&\qquad\longmapsto\mathscr Q_A
\longmapsto\mathfrak R_A^{\mathrm{all}}
\longmapsto\{F_R\}
\longmapsto V_A,r_A\\
&\qquad\longmapsto(\psi_1,\ldots,\psi_{r_A}),\varphi_A.
\end{aligned}
\]
No \(\mathcal D\), target, validity constraint, selected rule, or realized
seed enters this chain.  The seed law, rather than a realized seed, enters
each \(F_R\).  Every later valid-rule family is a subset of the already fixed
all-rule family.

If \(r_A=0\), then \(f=0\), \(a(f)\), and \(\varphi_A(x)\) are the empty
vectors in \(\mathbb R^0\), and their inner product is \(0=f(x)\).  Empty
\(\mathcal X\) forces this branch and has no pointwise condition.  Empty
\(\mathcal H\) leaves the construction unchanged and makes its later
target-indexed use vacuous.  At \(m=0\), the same basis is built from
\(\operatorname{span}\{F_{R_\varnothing}\}\), with no query-round argument
and with \(r_A\le B\).

Together with Lemma~\ref{lem:step-001-complete-rules} and
Proposition~\ref{prop:step-001-mean-response-space}, this constructs the
complete all-rule family, the exact average-before-span response space and
rank, and one common coordinate map before every learning instance.
\end{proof}

\subsection{Compactness in One Fixed Function Space}
\label{app:step-002}

\begin{lemma}[Evaluation coordinates determine the fixed topology]
\label{lem:step-002-evaluation-coordinates}
Under Assumptions~\ref{assump:adaptive-sq-interface} and
\ref{assump:mean-response-rank}, and the conclusions of
Propositions~\ref{prop:step-001-mean-response-space} and
\ref{prop:step-001-fixed-coordinates}, suppose \(r_A>0\).  Then
\[
\operatorname{span}_{\mathbb R}\{\delta_x:x\in\mathcal X\}=V_A^*,
\qquad \delta_x(f):=f(x).
\]
Consequently, there are points \(z_1,\ldots,z_{r_A}\in\mathcal X\) such that
\(\delta_{z_1},\ldots,\delta_{z_{r_A}}\) form a basis of \(V_A^*\).  The map
\[
E_Z:V_A\to\mathbb R^{r_A},
\qquad E_Z(f):=(f(z_1),\ldots,f(z_{r_A}))
\]
is a linear isomorphism, and
\(\|f\|_Z:=\|E_Z(f)\|_\infty\) is a norm inducing exactly the
finite-dimensional topology in which \(K_A\) was defined.
\end{lemma}

\begin{proof}
Every \(\delta_x\) is an algebraic linear functional on \(V_A\).  Put
\[
W:=\operatorname{span}_{\mathbb R}\{\delta_x:x\in\mathcal X\}
\subseteq V_A^*.
\]
The basis \(\psi_1,\ldots,\psi_{r_A}\) from
Proposition~\ref{prop:step-001-fixed-coordinates} shows directly that
\(\dim V_A^*=r_A\): the functionals returning the coefficients in
\(f=\sum_ja(f)_j\psi_j\) form a basis of the dual.

Suppose that \(W\ne V_A^*\), and write \(s=\dim W<r_A\).  Choose a basis
\(\lambda_1,\ldots,\lambda_s\) of \(W\).  In the fixed basis of \(V_A\),
the equations
\[
\lambda_i\!\left(\sum_{j=1}^{r_A}c_j\psi_j\right)=0,
\qquad 1\le i\le s,
\]
form an \(s\)-by-\(r_A\) homogeneous real linear system.  Row reduction has
at most \(s<r_A\) pivots, so the nullspace contains a nonzero vector \(c\).
Thus
\[
f:=\sum_{j=1}^{r_A}c_j\psi_j\ne0,
\qquad \lambda_i(f)=0\quad(1\le i\le s).
\]
Every \(\delta_x\) is a linear combination of the \(\lambda_i\), so
\(f(x)=\delta_x(f)=0\) for every \(x\in\mathcal X\).  Equality in
\(\mathbb R^{\mathcal X}\) is pointwise, so this makes \(f=0\), a
contradiction.  Hence \(W=V_A^*\).

Starting with the empty list, whenever selected evaluations do not yet span
\(V_A^*\), choose another evaluation outside their span.  Each addition
raises the dimension by one, so after exactly \(r_A\) additions the selected
evaluations form a basis.  Denote their points by
\(z_1,\ldots,z_{r_A}\).  In particular, \(r_A>0\) forces
\(\mathcal X\ne\varnothing\).

The map \(E_Z\) is linear.  If \(E_Z(f)=0\), every basis functional
\(\delta_{z_i}\) vanishes on \(f\).  Every \(\delta_x\) is their linear
combination, so \(f(x)=0\) for all \(x\), hence \(f=0\).  Thus \(E_Z\) is
injective.  Its domain and codomain both have dimension \(r_A\), so it is
surjective as well.

It follows that \(\|f\|_Z=\|E_Zf\|_\infty\) is a norm.  To identify its
topology with the fixed topology, let \(a(f)\) be the coordinate vector in
the basis \((\psi_j)\), and set
\[
M_{ij}:=\psi_j(z_i),\qquad 1\le i,j\le r_A.
\]
Then \(E_Z(f)=Ma(f)\), and bijectivity of \(E_Z\) makes \(M\) invertible.
Finite matrix multiplication gives
\[
\|E_Z(f)\|_\infty
\le
\left(\max_{1\le i\le r_A}\sum_{j=1}^{r_A}|M_{ij}|\right)
\|a(f)\|_\infty
\]
and
\[
\|a(f)\|_\infty
\le
\left(\max_{1\le i\le r_A}\sum_{j=1}^{r_A}|(M^{-1})_{ij}|\right)
\|E_Z(f)\|_\infty.
\]
Both multipliers are finite sums of fixed real entries.  Hence the evaluation
norm and the fixed-basis coordinate norm have the same open and closed sets.
The latter is the finite-dimensional topology specified in the setup, so
\(\|\cdot\|_Z\) induces that same topology.  These inequalities establish
topological equivalence only; they assume no bound on the coefficients of
\(F_R\) or of convex combinations.
\end{proof}

\begin{lemma}[Heine--Borel in evaluation coordinates]
\label{lem:step-002-heine-borel}
Under Assumptions~\ref{assump:adaptive-sq-interface} and
\ref{assump:mean-response-rank}, the conclusions of
Propositions~\ref{prop:step-001-mean-response-space} and
\ref{prop:step-001-fixed-coordinates}, and
Lemma~\ref{lem:step-002-evaluation-coordinates}, suppose \(r_A>0\).  If
\(C\subseteq V_A\) is closed in the fixed finite-dimensional topology and
bounded in \(\|\cdot\|_Z\), then \(C\) is compact in that topology.
\end{lemma}

\begin{proof}
Lemma~\ref{lem:step-002-evaluation-coordinates} makes \(E_Z\) an isometric
linear homeomorphism from \((V_A,\|\cdot\|_Z)\) onto
\((\mathbb R^{r_A},\|\cdot\|_\infty)\).  Hence \(D:=E_Z(C)\) is closed and
bounded in \(\mathbb R^{r_A}\).  We prove compactness directly.

Every sequence in \(D\) has a convergent subsequence with limit in \(D\).
Boundedness places every coordinate sequence in an interval \([-Q,Q]\).  A
bounded real sequence has a convergent subsequence by bisection: bisect the
interval, retain a closed half containing infinitely many terms, and iterate,
choosing successively increasing indices from the retained halves.  The
nested intervals have lengths tending to zero, so completeness of
\(\mathbb R\) gives their unique common limit and convergence of the chosen
subsequence.  Apply this argument successively to coordinates
\(1,\ldots,r_A\), each time taking a subsequence of the preceding one.  The
final subsequence converges in every coordinate and therefore in
\(\|\cdot\|_\infty\).  Closedness of \(D\) keeps the limit in \(D\).

This subsequence property first implies total boundedness.  Otherwise, for
some \(\eta>0\), one could recursively choose points whose pairwise distances
are at least \(\eta\); that sequence has no convergent subsequence.  Now let
\(\mathcal U\) be an open cover of \(D\).  If no \(\delta>0\) had the
property that each relative ball \(B_D(y,\delta)\) lies in some member of
\(\mathcal U\), one could choose \(y_n\in D\) such that
\(B_D(y_n,1/n)\) lies in no cover member.  A convergent subsequence
\(y_{n_j}\to y\), together with one cover member containing \(y\),
contradicts this choice for all sufficiently large \(j\).  Such a
\(\delta\) therefore exists.  A finite \(\delta/2\)-net gives finitely many
\(\delta\)-balls, each contained in a member of the cover, and hence a finite
subcover.  Thus \(D\) is compact.

The continuous inverse \(E_Z^{-1}\) maps \(D\) onto \(C\), proving that
\(C\) is compact in the fixed topology of \(V_A\).
\end{proof}

\begin{proposition}[The fixed response body is compact and convex]
\label{prop:step-002-compact-body}
Under Assumptions~\ref{assump:adaptive-sq-interface} and
\ref{assump:mean-response-rank}, the conclusions of
Lemma~\ref{lem:step-001-complete-rules} and
Proposition~\ref{prop:step-001-mean-response-space}, and
Lemmas~\ref{lem:step-002-evaluation-coordinates} and
\ref{lem:step-002-heine-borel}, suppose \(r_A>0\).  Then
\[
K_A=\overline{\operatorname{conv}}
\{F_R:R\in\mathfrak R_A^{\mathrm{all}}\}
\]
is a nonempty compact convex subset of \(V_A\) in the one fixed topology, and
\[
K_A\subseteq\{f\in V_A:\|f\|_Z\le1\}.
\]
\end{proposition}

\begin{proof}
Lemma~\ref{lem:step-001-complete-rules} makes the generating family nonempty.
For every complete rule \(R\) and every selected point \(z_i\),
Proposition~\ref{prop:step-001-mean-response-space} gives
\(|F_R(z_i)|\le1\).  Hence
\[
\|F_R\|_Z=\max_{1\le i\le r_A}|F_R(z_i)|\le1.
\]
If \(f=\sum_{j=1}^N\alpha_jF_{R_j}\) is a finite convex combination, then
for every \(i\),
\[
|f(z_i)|
\le\sum_{j=1}^N\alpha_j|F_{R_j}(z_i)|
\le\sum_{j=1}^N\alpha_j=1.
\]
Thus the convex hull lies in the evaluation-norm unit ball.  The reverse
triangle inequality
\[
|\|f\|_Z-\|g\|_Z|\le\|f-g\|_Z
\]
makes this ball closed.  Since
Lemma~\ref{lem:step-002-evaluation-coordinates} identifies its topology with
the topology in the defining closure, \(K_A\) remains in that unit ball.

By definition, \(K_A\) is closed, and it is nonempty because it contains the
nonempty convex hull.  To verify convexity of the closure, let
\(f,g\in K_A\), \(\theta\in[0,1]\), and \(\eta>0\).  Choose convex-hull
elements \(f_0,g_0\) satisfying
\(\|f-f_0\|_Z<\eta\) and \(\|g-g_0\|_Z<\eta\).  Then
\[
\|\theta f+(1-\theta)g-[\theta f_0+(1-\theta)g_0]\|_Z
\le\theta\|f-f_0\|_Z+(1-\theta)\|g-g_0\|_Z
<\eta.
\]
Since \(\theta f_0+(1-\theta)g_0\) belongs to the convex hull and
\(\eta\) is arbitrary, \(\theta f+(1-\theta)g\in K_A\).  Finally,
Lemma~\ref{lem:step-002-heine-borel} turns the established closedness and
boundedness into compactness in the same topology.
\end{proof}

\begin{lemma}[Every point evaluation is continuous with exact reconstruction]
\label{lem:step-002-continuous-evaluations}
Under Assumptions~\ref{assump:adaptive-sq-interface} and
\ref{assump:mean-response-rank}, the conclusions of
Propositions~\ref{prop:step-001-mean-response-space} and
\ref{prop:step-001-fixed-coordinates}, and
Lemma~\ref{lem:step-002-evaluation-coordinates}, suppose \(r_A>0\).  For
every \(x\in\mathcal X\), point evaluation \(\delta_x(f)=f(x)\) is
continuous in the fixed topology.  If \(e_1,\ldots,e_{r_A}\) are the
standard coordinate vectors and
\[
b_i(x):=(E_Z^{-1}e_i)(x),
\]
then every \(f\in V_A\) satisfies
\[
f(x)=\sum_{i=1}^{r_A}b_i(x)f(z_i),
\qquad
|f(x)|\le\left(\sum_{i=1}^{r_A}|b_i(x)|\right)\|f\|_Z,
\]
and
\[
\delta_x(f)-(\delta_x\circ E_Z^{-1})(E_Zf)=0.
\]
\end{lemma}

\begin{proof}
Fix \(x\).  The composition \(\ell_x:=\delta_x\circ E_Z^{-1}\) is linear
on \(\mathbb R^{r_A}\).  For
\(y=\sum_{i=1}^{r_A}y_ie_i\), linearity gives
\[
\ell_x(y)=\sum_{i=1}^{r_A}y_i\ell_x(e_i)
=\sum_{i=1}^{r_A}y_i b_i(x).
\]
Therefore
\[
|\ell_x(y)|
\le\sum_{i=1}^{r_A}|b_i(x)|\,|y_i|
\le\left(\sum_{i=1}^{r_A}|b_i(x)|\right)\|y\|_\infty.
\]
The coefficient sum is finite for each fixed \(x\), proving continuity.
Substituting \(y=E_Zf\) and using \(E_Z^{-1}E_Zf=f\) yields both the
reconstruction formula and the zero-residual factorization.  No bound uniform
in \(x\) is asserted or needed.
\end{proof}

\begin{proposition}[The zero-dimensional and empty-domain branch]
\label{prop:step-002-zero-dimensional}
Under Assumptions~\ref{assump:adaptive-sq-interface} and
\ref{assump:mean-response-rank}, and the conclusions of
Lemma~\ref{lem:step-001-complete-rules} and
Proposition~\ref{prop:step-001-mean-response-space}, suppose \(r_A=0\).
Then
\[
V_A=K_A=\{0\}.
\]
The empty evaluation map \(E_\varnothing:V_A\to\mathbb R^0\) is a linear
isomorphism and induces the unique zero-dimensional topology.  The body
\(K_A\) is nonempty, compact, and convex, and every point evaluation is the
zero continuous functional.  These conclusions include
\(\mathcal X=\varnothing\) and the zero-rank branch at \(m=0\).
\end{proposition}

\begin{proof}
A vector space of dimension zero contains only its zero vector, so
\(V_A=\{0\}\).  Lemma~\ref{lem:step-001-complete-rules} makes the generator
family nonempty, and every generator lies in \(V_A\); hence every generator
is zero.  Its convex hull and its closure are both \(\{0\}\), proving
\(K_A=\{0\}\).

Both \(V_A\) and \(\mathbb R^0\) are singleton vector spaces, so the empty
evaluation map is directly linear, injective, and surjective.  The singleton
\(\{0\}\) is convex and compact by the open-cover definition.  Every point
evaluation has domain \(\{0\}\) and value zero, so it is continuous.  If
\(\mathcal X=\varnothing\), there is no point evaluation to check.  If
\(m=0\), the unique empty response rule supplies the generator without any
query or reply.
\end{proof}

\begin{proof}[Consequences for the fixed response body]
If \(r_A>0\), Lemma~\ref{lem:step-002-evaluation-coordinates} supplies actual
domain evaluations and identifies their norm with the fixed topology;
Proposition~\ref{prop:step-002-compact-body} gives a nonempty compact convex
\(K_A\); and Lemma~\ref{lem:step-002-continuous-evaluations} makes every
point evaluation continuous with exact reconstruction.  If \(r_A=0\),
Proposition~\ref{prop:step-002-zero-dimensional} proves all the same body and
continuity conclusions directly.  Thus the later minimax and intersection
arguments always take place in one fixed compact body, including for empty
\(\mathcal X\) and \(m=0\).
\end{proof}

\subsection{Exact-Center Correlation and Nondegeneracy}
\label{app:step-003}

\begin{lemma}[Exact-center legality]
\label{lem:step-003-exact-center-legality}
Under Assumptions~\ref{assump:parameter-regime} and
\ref{assump:adaptive-sq-interface}, and
Lemma~\ref{lem:step-001-complete-rules}, for every
\(\mathcal D\in\mathcal P(\mathcal X)\) and \(h\in\mathcal H\), define on
every admitted public query-bearing history
\[
R^0_{\mathcal D,h}(q_1,v_1,\ldots,q_t)
:=\mathbb E_{x\sim\mathcal D}q_t(x,h(x)).
\]
Then \(R^0_{\mathcal D,h}\) is a deterministic complete public-transcript
rule in \(\mathfrak R_A^{\mathrm{all}}\) and belongs to
\(\mathfrak R_{A,\tau}(\mathcal D,h)\).  Its validity deviation is zero at
every round reached by every seed.  For \(m=0\), it is the unique empty rule
and validity is vacuous.
\end{lemma}

\begin{proof}
Fix \(\mathcal D\) and \(h\).  At any admitted public query-bearing history,
the current query is measurable, bounded, and \([-1,1]\)-valued.  Therefore
its population expectation is well defined and obeys
\[
-1\le\mathbb E_{x\sim\mathcal D}q_t(x,h(x))\le1.
\]
Since \(\tau>0\), \([-1,1]\subseteq I_\tau\), so the displayed formula
assigns an admissible reply to every admitted history, including histories
not reached by one seed or preceding reply path.

For fixed \((\mathcal D,h)\), the formula is deterministic and takes as its
runtime input only the public transcript, in particular the public query
\(q_t\); it has no hidden-seed input.  Lemma~\ref{lem:step-001-complete-rules}
therefore identifies it as a complete rule in
\(\mathfrak R_A^{\mathrm{all}}\).

On every reached seed path and at every reached round, its reply is the
population center itself, so
\[
\left|R^0_{\mathcal D,h}(q_1,v_1,\ldots,q_t)
-\mathbb E_{x\sim\mathcal D}q_t(x,h(x))\right|=0\le\tau.
\]
This proves validity separately on every reached path.  There is no
accumulated tolerance error over the at most \(m\) rounds.  For \(m=0\),
there is no query-bearing history, and
Lemma~\ref{lem:step-001-complete-rules} gives the unique empty rule.
\end{proof}

\begin{proposition}[Exact seed-averaged correlation]
\label{prop:step-003-exact-correlation}
Under Assumptions~\ref{assump:parameter-regime},
\ref{assump:adaptive-sq-interface}, and
\ref{assump:universal-adversarial-guarantee}, and the conclusions of
Proposition~\ref{prop:step-001-mean-response-space} and
Lemma~\ref{lem:step-003-exact-center-legality}, for every
\(\mathcal D\in\mathcal P(\mathcal X)\) and \(h\in\mathcal H\),
\[
F_{R^0_{\mathcal D,h}}
\in\{F_R:R\in\mathfrak R_A^{\mathrm{all}}\}
\subseteq K_A\subseteq V_A
\]
and
\[
\begin{aligned}
\mathbb E_{x\sim\mathcal D}[h(x)F_{R^0_{\mathcal D,h}}(x)]
&=1-2\,\mathbb E_{U\sim\mu_A}
\mathcal L_{\mathcal D,h}(g_{U,R^0_{\mathcal D,h}})\\
&\ge1-2\varepsilon=\rho>\frac12.
\end{aligned}
\]
The accuracy expectation is solely over the learner seed; no distribution
over response policies and no favorable seed is introduced.
\end{proposition}

\begin{proof}
Lemma~\ref{lem:step-003-exact-center-legality} gives
\(R^0_{\mathcal D,h}\in\mathfrak R_A^{\mathrm{all}}\).  Hence
Proposition~\ref{prop:step-001-mean-response-space} defines its mean response
and places it in the raw generating family.  Every raw generator belongs to
its convex hull and therefore to \(K_A\subseteq V_A\); compactness is not
needed for this inclusion.

The bounded Fubini theorem \citep{folland1999real} applies with probability
laws \(\mu_A\) and \(\mathcal D\), and with
\[
a(u,x):=h(x)g_{u,R^0_{\mathcal D,h}}(x).
\]
Assumption~\ref{assump:adaptive-sq-interface} and the setup supply the needed
measurability, while \(|a(u,x)|=1\) gives integrability.  Thus
\[
\begin{aligned}
\mathbb E_{x\sim\mathcal D}[h(x)F_{R^0_{\mathcal D,h}}(x)]
&=\int_{\mathcal X}h(x)
\left(\int_{\Omega_A}g_{u,R^0_{\mathcal D,h}}(x)\,\mu_A(du)\right)
\mathcal D(dx)\\
&=\int_{\Omega_A}\left(\int_{\mathcal X}
h(x)g_{u,R^0_{\mathcal D,h}}(x)\,\mathcal D(dx)\right)\mu_A(du).
\end{aligned}
\]
For every \((u,x)\), both binary labels lie in \(\{-1,+1\}\), and the two
possible cases give the pointwise identity
\[
h(x)g_{u,R^0_{\mathcal D,h}}(x)
=1-2\mathbf 1\{g_{u,R^0_{\mathcal D,h}}(x)\ne h(x)\}.
\]
Integration over \(x\) gives
\[
\int_{\mathcal X}h(x)g_{u,R^0_{\mathcal D,h}}(x)\,\mathcal D(dx)
=1-2\mathcal L_{\mathcal D,h}(g_{u,R^0_{\mathcal D,h}}),
\]
and integration over the seed yields the exact identity in the statement.

Lemma~\ref{lem:step-003-exact-center-legality} proves that this particular
deterministic complete rule is valid.  The universal quantifier in
Assumption~\ref{assump:universal-adversarial-guarantee} therefore applies to
it and gives
\[
\mathbb E_{U\sim\mu_A}
\mathcal L_{\mathcal D,h}(g_{U,R^0_{\mathcal D,h}})\le\varepsilon.
\]
Substitution, with no discarded term, gives
\[
1-2\mathbb E_U\mathcal L_{\mathcal D,h}(g_{U,R^0})
\ge1-2\varepsilon=\rho>\frac12.
\]
The exact-center rule is only one rule already covered by the every-valid-rule
premise; it does not replace the adversarial oracle model.  Tolerance creates
zero residual.  At \(m=0\), the same calculation uses the empty rule and no
query-round step.  At \(\varepsilon=0\), nonnegativity of loss forces its seed
expectation to be zero, so the correlation is exactly \(1\).
\end{proof}

\begin{proposition}[Point-mass exclusion of zero rank]
\label{prop:step-003-nonzero-rank}
Under Assumptions~\ref{assump:parameter-regime},
\ref{assump:adaptive-sq-interface}, and
\ref{assump:universal-adversarial-guarantee}, and the conclusions of
Propositions~\ref{prop:step-001-mean-response-space} and
\ref{prop:step-003-exact-correlation}, if
\(\mathcal X\ne\varnothing\) and \(\mathcal H\ne\varnothing\), then for
every \(x_0\in\mathcal X\) and \(h\in\mathcal H\),
\[
h(x_0)F_{R^0_{\delta_{x_0},h}}(x_0)\ge\rho>0.
\]
Consequently \(V_A\ne\{0\}\) and \(r_A\ge1\).
\end{proposition}

\begin{proof}
For nonempty \(\mathcal X\), the point mass \(\delta_{x_0}\) belongs to
\(\mathcal P(\mathcal X)\).  Proposition~\ref{prop:step-003-exact-correlation}
with \(\mathcal D=\delta_{x_0}\) gives
\[
\begin{aligned}
h(x_0)F_{R^0_{\delta_{x_0},h}}(x_0)
&=\mathbb E_{x\sim\delta_{x_0}}
[h(x)F_{R^0_{\delta_{x_0},h}}(x)]\\
&\ge\rho>\frac12>0.
\end{aligned}
\]
Thus this raw response is a nonzero function in \(V_A\), and
\(r_A=\dim V_A\ge1\).  When \(m=0\), all these exact-center rules are the
unique empty rule, but the point-mass conclusion is unchanged.  If
\(\mathcal X\) or \(\mathcal H\) is empty, no nonexistent point or target is
selected, and the later sign condition is vacuous.

Together with Lemma~\ref{lem:step-003-exact-center-legality} and
Proposition~\ref{prop:step-003-exact-correlation}, this proves that every
distribution-target pair has an exact \(\rho\)-correlated member of the one
fixed response body and that the zero-rank branch is impossible precisely in
the nonempty-domain, nonempty-class regime.
\end{proof}

\subsection{Finite Simultaneous Margin by Minimax}
\label{app:step-004}

\begin{lemma}[Finite-simplex minima occur at vertices]
\label{lem:step-004-simplex-vertex}
Let \(S\subseteq\mathcal X\) be nonempty and finite.  For every real family
\((a_x)_{x\in S}\),
\[
\min_{p\in\Delta(S)}\sum_{x\in S}p(x)a_x
=\min_{x\in S}a_x,
\]
where
\[
\Delta(S):=\left\{p\in\mathbb R^S:p(x)\ge0\ \text{for all }x\in S,
\ \sum_{x\in S}p(x)=1\right\}.
\]
Both minima are attained at a simplex vertex, including when \(S\) is a
singleton.
\end{lemma}

\begin{proof}
Choose \(x_*\in S\) with \(a_{x_*}=\min_{x\in S}a_x\).  For every
\(p\in\Delta(S)\),
\[
\sum_{x\in S}p(x)a_x
\ge\sum_{x\in S}p(x)a_{x_*}
=a_{x_*}.
\]
At the vertex \(p^{x_*}\), which assigns mass one to \(x_*\), equality
holds.  This calculation also covers \(S=\{x_*\}\), for which the simplex
contains only that vertex.
\end{proof}

\begin{proposition}[Sion's equality on the fixed response body]
\label{prop:step-004-sion-fixed-body}
Under Assumptions~\ref{assump:adaptive-sq-interface} and
\ref{assump:mean-response-rank}, and the conclusions of
Propositions~\ref{prop:step-002-compact-body} and
\ref{prop:step-002-zero-dimensional},
Lemma~\ref{lem:step-002-continuous-evaluations}, and
Lemma~\ref{lem:step-004-simplex-vertex}, let \(h\in\mathcal H\) and let
\(S\subseteq\mathcal X\) be nonempty and finite.  On the fixed product
\(K_A\times\Delta(S)\), define
\[
L_h(f,p):=\sum_{x\in S}p(x)h(x)f(x).
\]
Then
\[
\max_{f\in K_A}\min_{p\in\Delta(S)}L_h(f,p)
=\min_{p\in\Delta(S)}\max_{f\in K_A}L_h(f,p),
\]
and every displayed maximum and minimum is attained.
\end{proposition}

\begin{proof}
We use Sion's Theorem 3.4 \citep[Theorem~3.4]{sion1958minimax} in its
maximizing/minimizing orientation.  The cited statement says that if \(M\)
and \(N\) are nonempty convex compact spaces contained in linear topological
spaces, and \(G:M\times N\to\mathbb R\) is upper semicontinuous and
quasiconcave in its first, maximizing variable for every fixed second
variable, and lower semicontinuous and quasiconvex in its second, minimizing
variable for every fixed first variable, then
\[
\sup_{m\in M}\inf_{n\in N}G(m,n)
=\inf_{n\in N}\sup_{m\in M}G(m,n).
\tag{A.1}
\]

Proposition~\ref{prop:step-002-compact-body} gives a nonempty compact convex
\(K_A\) when \(r_A>0\), and
Proposition~\ref{prop:step-002-zero-dimensional} gives the same conclusion
directly when \(r_A=0\).  Thus the maximizing set is always the same fixed
body; it is not selected after \(h\), \(S\), or \(p\).

Because \(S\) is nonempty and finite, \(\Delta(S)\) contains every vertex,
is convex, and is closed and bounded in the finite-dimensional linear
topological space \(\mathbb R^S\).  The direct proof of
Lemma~\ref{lem:step-002-heine-borel} applies with the exact substitutions
\(r_A\mapsto |S|\), \(D\mapsto\Delta(S)\), and the selected evaluation
coordinates replaced by the standard coordinates of \(\mathbb R^S\).
Coordinatewise bisection gives a convergent subsequence, closedness retains
its limit, and the same finite-net and open-cover argument yields a finite
subcover.  Hence \(\Delta(S)\) is compact.  Thus both sets required by the
cited theorem are nonempty, compact, and convex.

For fixed \(p\in\Delta(S)\),
\[
f\longmapsto L_h(f,p)
=\sum_{x\in S}p(x)h(x)\,\delta_x(f)
\]
is a finite linear combination of point evaluations.
Lemma~\ref{lem:step-002-continuous-evaluations} gives their continuity in the
positive-rank branch, and
Proposition~\ref{prop:step-002-zero-dimensional} gives it in the zero-rank
branch.  Hence this map is continuous affine, therefore upper semicontinuous
and quasiconcave, on \(K_A\).  For fixed \(f\in K_A\),
\[
p\longmapsto L_h(f,p)
=\sum_{x\in S}p(x)h(x)f(x)
\]
is a finite linear function of the Euclidean coordinates of \(p\).  It is
continuous affine, hence lower semicontinuous and quasiconvex, on
\(\Delta(S)\).  The payoff is finite and real valued.

The mapping
\[
M=K_A,\qquad N=\Delta(S),\qquad G=L_h
\]
therefore discharges every hypothesis of Sion's theorem.  Its conclusion,
without reversing the players or changing signs, is
\[
\sup_{f\in K_A}\inf_{p\in\Delta(S)}L_h(f,p)
=\inf_{p\in\Delta(S)}\sup_{f\in K_A}L_h(f,p).
\tag{A.2}
\]
The cited theorem supplies this order equality only; the fixed body,
per-distribution correlation, vertex reduction, attainment, and later
arbitrary-domain closure are established separately in this paper.

We now justify all four attainment assertions.  For fixed \(p\), continuity
of \(L_h(\cdot,p)\) on compact \(K_A\) gives a maximum.  For fixed \(f\),
continuity of \(L_h(f,\cdot)\) on compact \(\Delta(S)\) gives a minimum;
Lemma~\ref{lem:step-004-simplex-vertex} identifies an attaining vertex.

For the outer maximum, that lemma also gives
\[
\inf_{p\in\Delta(S)}L_h(f,p)
=\min_{x\in S}h(x)f(x).
\]
This function is continuous in \(f\), because for \(f,g\in K_A\),
\[
\left|\min_{x\in S}h(x)f(x)-\min_{x\in S}h(x)g(x)\right|
\le\max_{x\in S}|f(x)-g(x)|,
\]
and the finitely many point evaluations are continuous.  Compactness of
\(K_A\) therefore makes the left outer supremum in (A.2) an attained maximum.

For the outer minimum, write
\(v(p):=\max_{f\in K_A}L_h(f,p)\).  Every quantity
\(\max_{f\in K_A}|f(x)|\) is finite because evaluation is continuous on
compact \(K_A\).  For \(p,q\in\Delta(S)\),
\[
\begin{aligned}
|v(p)-v(q)|
&\le\max_{f\in K_A}|L_h(f,p)-L_h(f,q)|\\
&\le\sum_{x\in S}|p(x)-q(x)|
\max_{f\in K_A}|f(x)|.
\end{aligned}
\]
Thus \(v\) is continuous on compact \(\Delta(S)\), so the right outer
infimum in (A.2) is attained.  Replacing all four suprema and infima by the
attained maxima and minima proves the proposition.
\end{proof}

\begin{proposition}[Exact-center lower bound for the fixed-body game]
\label{prop:step-004-exact-center-lower}
Under Assumptions~\ref{assump:parameter-regime},
\ref{assump:adaptive-sq-interface},
\ref{assump:universal-adversarial-guarantee}, and
\ref{assump:mean-response-rank}, and the conclusions of
Propositions~\ref{prop:step-003-exact-correlation} and
\ref{prop:step-004-sion-fixed-body}, let \(h\in\mathcal H\) and let
\(S\subseteq\mathcal X\) be nonempty and finite.  Then
\[
\min_{p\in\Delta(S)}\max_{f\in K_A}L_h(f,p)\ge\rho.
\]
For each \(p\), the exact-center rule and its mean response may depend on
\(p\), but every such response belongs to the one fixed \(K_A\).
\end{proposition}

\begin{proof}
Fix \(p\in\Delta(S)\).  It defines the finitely supported law
\(\mathcal D_p\in\mathcal P(\mathcal X)\) assigning mass \(p(x)\) to
\(x\in S\) and zero elsewhere.  Proposition~\ref{prop:step-003-exact-correlation}
with \((\mathcal D_p,h)\) gives
\(F_{R^0_{\mathcal D_p,h}}\in K_A\) and
\[
\begin{aligned}
L_h(F_{R^0_{\mathcal D_p,h}},p)
&=\sum_{x\in S}p(x)h(x)F_{R^0_{\mathcal D_p,h}}(x)\\
&=\mathbb E_{x\sim\mathcal D_p}
[h(x)F_{R^0_{\mathcal D_p,h}}(x)]\\
&\ge\rho.
\end{aligned}
\]
Consequently,
\[
\max_{f\in K_A}L_h(f,p)
\ge L_h(F_{R^0_{\mathcal D_p,h}},p)
\ge\rho.
\tag{A.3}
\]
The quantifiers are
\[
\forall p\in\Delta(S)\ \exists F_{R^0_{\mathcal D_p,h}}\in K_A:
L_h(F_{R^0_{\mathcal D_p,h}},p)\ge\rho.
\]
The witness may change with \(p\), but the maximizing set does not: it is
always the pre-instance body \(K_A\).  Taking the attained minimum of (A.3)
over \(p\) preserves \(\rho\) and proves the claim.
\end{proof}

\begin{proposition}[Attained simultaneous margin on every finite restriction]
\label{prop:step-004-finite-margin}
Under Assumptions~\ref{assump:parameter-regime}--\ref{assump:mean-response-rank},
and the conclusions of Lemma~\ref{lem:step-004-simplex-vertex} and
Propositions~\ref{prop:step-004-sion-fixed-body} and
\ref{prop:step-004-exact-center-lower}, for every \(h\in\mathcal H\) and
every nonempty finite \(S\subseteq\mathcal X\),
\[
\begin{aligned}
\max_{f\in K_A}\min_{x\in S}h(x)f(x)
&=\max_{f\in K_A}\min_{p\in\Delta(S)}L_h(f,p)\\
&=\min_{p\in\Delta(S)}\max_{f\in K_A}L_h(f,p)\\
&\ge\rho.
\end{aligned}
\]
There is an attaining \(f_{h,S}\in K_A\) satisfying
\[
h(x)f_{h,S}(x)\ge\rho
\qquad\text{for every }x\in S.
\]
\end{proposition}

\begin{proof}
For fixed \(f\in K_A\), apply
Lemma~\ref{lem:step-004-simplex-vertex} to
\(a_x=h(x)f(x)\).  This gives
\[
\min_{p\in\Delta(S)}L_h(f,p)
=\min_{x\in S}h(x)f(x).
\]
Taking the attained maximum over the same fixed \(K_A\) proves the first
equality.  Proposition~\ref{prop:step-004-sion-fixed-body} proves the second
equality in the displayed orientation, and
Proposition~\ref{prop:step-004-exact-center-lower} proves the final inequality
at the unchanged threshold \(\rho\).

The maximum is attained by
Proposition~\ref{prop:step-004-sion-fixed-body}; call an attaining function
\(f_{h,S}\).  Then
\[
\min_{x\in S}h(x)f_{h,S}(x)\ge\rho,
\]
which is equivalent, because \(S\ne\varnothing\), to all the pointwise
inequalities in the statement.

For \(S=\{x_0\}\), the simplex contains the single point mass and the chain
reduces to
\[
\max_{f\in K_A}h(x_0)f(x_0)
\ge h(x_0)F_{R^0_{\delta_{x_0},h}}(x_0)
\ge\rho.
\]
The empty set is not passed to Sion's theorem; its constraint is handled as
the empty intersection in Proposition~\ref{prop:step-005-finite-intersections}
below.  At \(\varepsilon=0\), all equalities preserve the margin \(1\).
\end{proof}

\subsection{Arbitrary-Domain Exactification}
\label{app:step-005}

\begin{lemma}[Closed signed constraints in the fixed body]
\label{lem:step-005-closed-constraints}
Under Assumptions~\ref{assump:adaptive-sq-interface} and
\ref{assump:mean-response-rank}, and the conclusions of
Lemma~\ref{lem:step-002-continuous-evaluations} and
Proposition~\ref{prop:step-002-zero-dimensional}, fix \(h\in\mathcal H\) and
\(x\in\mathcal X\).  Then
\[
C_{h,x}:=\{f\in K_A:h(x)f(x)\ge\rho\}
\]
is relatively closed in the one fixed body \(K_A\), at the same threshold as
Proposition~\ref{prop:step-004-finite-margin}, including \(\rho=1\).
\end{lemma}

\begin{proof}
Let \(\delta_x(f)=f(x)\).  In the positive-rank branch,
Lemma~\ref{lem:step-002-continuous-evaluations} makes \(\delta_x\) continuous
in the fixed topology of \(V_A\).  In the zero-rank branch,
Proposition~\ref{prop:step-002-zero-dimensional} makes it the zero continuous
functional.  Thus
\[
T_{h,x}:K_A\to\mathbb R,
\qquad T_{h,x}(f):=h(x)\delta_x(f),
\]
is continuous in every rank branch.  The ray \([\rho,\infty)\) is closed in
\(\mathbb R\), and
\[
C_{h,x}=T_{h,x}^{-1}([\rho,\infty)).
\]
Therefore \(C_{h,x}\) is closed in the subspace topology of the same
\(K_A\).  The weak inequality and the numerical threshold are unchanged at
\(\rho=1\).
\end{proof}

\begin{proposition}[Finite feasibility in the fixed body]
\label{prop:step-005-finite-intersections}
Assume Assumptions~\ref{assump:parameter-regime},
\ref{assump:adaptive-sq-interface},
\ref{assump:universal-adversarial-guarantee}, and
\ref{assump:mean-response-rank}.  Under the nonemptiness conclusions of
Propositions~\ref{prop:step-002-compact-body} and
\ref{prop:step-002-zero-dimensional}, together with
Proposition~\ref{prop:step-004-finite-margin}, fix \(h\in\mathcal H\).  Then
for every finite \(S\subseteq\mathcal X\), including \(S=\varnothing\),
\[
\bigcap_{x\in S}C_{h,x}\ne\varnothing.
\]
Thus \((C_{h,x})_{x\in\mathcal X}\) has the finite-intersection property
inside one fixed \(K_A\).
\end{proposition}

\begin{proof}
For \(S=\varnothing\), the intersection over the empty subfamily equals its
ambient set:
\[
\bigcap_{x\in\varnothing}C_{h,x}=K_A.
\]
Proposition~\ref{prop:step-002-compact-body} in the positive-rank branch and
Proposition~\ref{prop:step-002-zero-dimensional} in the zero-rank branch make
this set nonempty.

Now suppose \(S\ne\varnothing\).  Proposition~\ref{prop:step-004-finite-margin}
supplies \(f_{h,S}\in K_A\) such that
\[
h(x)f_{h,S}(x)\ge\rho
\qquad\text{for every }x\in S.
\]
By definition, \(f_{h,S}\in C_{h,x}\) for every \(x\in S\), and therefore
\[
f_{h,S}\in\bigcap_{x\in S}C_{h,x}.
\]
These two cases exhaust all finite subfamilies.  The witnesses for different
finite sets need not agree, be nested, or converge.
\end{proof}

\begin{lemma}[Compactness closes an arbitrary finite-intersection family]
\label{lem:step-005-compact-fip}
Let \(K\) be a nonempty compact topological space, let \(I\) be an arbitrary,
possibly empty, index set, and let \((D_i)_{i\in I}\) be relatively closed
subsets of \(K\).  If
\[
\bigcap_{i\in J}D_i\ne\varnothing
\qquad\text{for every finite }J\subseteq I,
\]
where the intersection for \(J=\varnothing\) is \(K\), then
\[
\bigcap_{i\in I}D_i\ne\varnothing.
\]
\end{lemma}

\begin{proof}
Suppose instead that \(\bigcap_{i\in I}D_i=\varnothing\).  Taking
complements relative to \(K\) gives
\[
K=K\setminus\bigcap_{i\in I}D_i
=\bigcup_{i\in I}(K\setminus D_i).
\]
Each \(K\setminus D_i\) is relatively open, so these complements form an
open cover of \(K\).  Compactness supplies finitely many indices
\(i_1,\ldots,i_n\) whose complements already cover \(K\).  Since
\(K\ne\varnothing\), an empty subcollection cannot cover it, so \(n\ge1\).
Taking complements again gives
\[
\bigcap_{j=1}^nD_{i_j}=\varnothing,
\]
contradicting finite-intersection nonemptiness.  Hence the total intersection
is nonempty.  The argument uses the open-cover definition of compactness and
does not enumerate \(I\), pass to a sequence, or impose a cardinality
condition.  If \(I=\varnothing\), the conclusion is simply
\(K\ne\varnothing\).
\end{proof}

\begin{proposition}[One exact-margin witness on the arbitrary domain]
\label{prop:step-005-global-margin}
Under Assumptions~\ref{assump:parameter-regime}--\ref{assump:mean-response-rank},
and the conclusions of Propositions~\ref{prop:step-002-compact-body},
\ref{prop:step-002-zero-dimensional},
\ref{prop:step-004-finite-margin}, and
\ref{prop:step-005-finite-intersections}, and
Lemmas~\ref{lem:step-002-continuous-evaluations},
\ref{lem:step-005-closed-constraints}, and
\ref{lem:step-005-compact-fip}, for every \(h\in\mathcal H\) there exists
\(f_h\in K_A\) such that
\[
h(x)f_h(x)\ge\rho
\qquad\text{for every }x\in\mathcal X.
\]
The conclusion holds for an arbitrary, possibly uncountable or empty, domain
and preserves \(\rho\) exactly, including \(\rho=1\).
\end{proposition}

\begin{proof}
Fix \(h\in\mathcal H\).  Proposition~\ref{prop:step-002-compact-body} in the
positive-rank branch and Proposition~\ref{prop:step-002-zero-dimensional} in
the zero-rank branch make the same \(K_A\) nonempty and compact.
Lemma~\ref{lem:step-005-closed-constraints} makes every \(C_{h,x}\) relatively
closed in that body, and
Proposition~\ref{prop:step-005-finite-intersections} gives the
finite-intersection property, including the empty subfamily.  Apply
Lemma~\ref{lem:step-005-compact-fip} with
\[
K=K_A,\qquad I=\mathcal X,\qquad D_x=C_{h,x}.
\]
It follows that
\[
\bigcap_{x\in\mathcal X}C_{h,x}\ne\varnothing.
\]
Choose \(f_h\) in this intersection.  For every \(x\), its membership in
\(C_{h,x}\) is exactly \(h(x)f_h(x)\ge\rho\).

No relationship among the finite witnesses \(f_{h,S}\) was used.  They may
vary arbitrarily with \(S\); their only role is to certify finite
intersections inside the same \(K_A\), while compactness produces \(f_h\).

If \(\mathcal X=\varnothing\), the indexed family is empty and its full
intersection is \(K_A\); the pointwise condition is vacuous.  If
\(\mathcal H=\varnothing\), the outer target quantifier is vacuous.  In the
zero-rank branch, \(K_A=\{0\}\).  Empty \(\mathcal X\) is compatible with
this branch.  If instead \(\mathcal X\ne\varnothing\) and some
\(h\in\mathcal H\) exists, Proposition~\ref{prop:step-004-finite-margin}
applied to \(S=\{x\}\) would require \(h(x)0\ge\rho>1/2\), a contradiction.
Thus the preceding results themselves exclude zero rank exactly in the
nonempty target/domain regime; no positive-rank assumption was added.

At \(\rho=1\), every finite witness belongs to the closed sets defined by
\(h(x)f(x)\ge1\), and the compact intersection uses those identical sets.
No relaxation, sequential approximation, or smaller positive threshold is
introduced.  This proves the arbitrary-domain conclusion at zero residual.
\end{proof}

\subsection{Exact Scores and Polynomial Dimension Closure}
\label{app:step-006}

\begin{proposition}[Exact score transfer through the pre-instance basis]
\label{prop:step-006-exact-score-transfer}
Under Assumptions~\ref{assump:parameter-regime}--\ref{assump:mean-response-rank},
and the conclusions of Propositions~\ref{prop:step-001-fixed-coordinates} and
\ref{prop:step-005-global-margin}, the basis
\(\psi_1,\ldots,\psi_{r_A}\) and map \(\varphi_A\) are fixed before every
learning instance.  For every \(h\in\mathcal H\) and every \(f_h\in K_A\)
supplied by Proposition~\ref{prop:step-005-global-margin}, there is a unique
\(w_h\in\mathbb R^{r_A}\) such that
\[
f_h=\sum_{j=1}^{r_A}(w_h)_j\psi_j
\]
and, for every \(x\in\mathcal X\),
\[
\langle w_h,\varphi_A(x)\rangle=f_h(x),
\qquad
h(x)\langle w_h,\varphi_A(x)\rangle
\ge\rho>\frac12>0.
\]
Thus every nonvacuous score is nonzero and has sign \(h(x)\).  Only
\(f_h\) and \(w_h\) may depend on \(h\).
\end{proposition}

\begin{proof}
Proposition~\ref{prop:step-001-fixed-coordinates} constructs the basis and
\(\varphi_A\) from the all-rule mean-response space, before any distribution,
target, validity restriction, selected rule, or realized seed.  Fixing a
later target cannot change this basis or map.

Now fix \(h\in\mathcal H\).  Proposition~\ref{prop:step-005-global-margin}
supplies \(f_h\in K_A\) with
\[
h(x)f_h(x)\ge\rho
\qquad\text{for every }x\in\mathcal X.
\]
The closure defining \(K_A\) is taken inside \(V_A\), so \(f_h\in V_A\).
The basis property gives a unique \(w_h\in\mathbb R^{r_A}\) satisfying
\[
f_h=\sum_{j=1}^{r_A}(w_h)_j\psi_j.
\]
Evaluating at any \(x\) gives
\[
\begin{aligned}
f_h(x)
&=\sum_{j=1}^{r_A}(w_h)_j\psi_j(x)\\
&=\left\langle w_h,
(\psi_1(x),\ldots,\psi_{r_A}(x))\right\rangle\\
&=\langle w_h,\varphi_A(x)\rangle.
\end{aligned}
\]
This is an equality, with no limit, norm conversion, approximation, or
coordinate residual.  Multiplying by \(h(x)\) transfers the global margin:
\[
h(x)\langle w_h,\varphi_A(x)\rangle
=h(x)f_h(x)\ge\rho.
\]
Assumption~\ref{assump:parameter-regime} gives \(\rho>1/2>0\).  Hence when
\(h(x)=+1\) the score is positive, and when \(h(x)=-1\) it is negative.

If \(\mathcal H=\varnothing\), the target statement is vacuous and the map
remains fixed.  If \(\mathcal X=\varnothing\), all pointwise assertions are
vacuous; Proposition~\ref{prop:step-001-mean-response-space} gives
\(V_A=\{0\}\) and \(r_A=0\), so \(f_h=0\), \(w_h\) is the empty vector, and
the basis expansion is the empty sum.
\end{proof}

\begin{lemma}[Dimension-complexity admissibility and degenerate branches]
\label{lem:step-006-dimension-admissibility}
Assume Assumptions~\ref{assump:parameter-regime},
\ref{assump:adaptive-sq-interface},
\ref{assump:universal-adversarial-guarantee}, and
\ref{assump:mean-response-rank}.  Under the conclusions of
Proposition~\ref{prop:step-006-exact-score-transfer}
and Proposition~\ref{prop:step-003-nonzero-rank},
\[
\operatorname{dc}(\mathcal H)\le r_A.
\]
If \(\mathcal X=\varnothing\) or \(\mathcal H=\varnothing\), then
\(\operatorname{dc}(\mathcal H)=0\).  If \(r_A=0\), at least one of these
two sets is empty and
\(\operatorname{dc}(\mathcal H)=r_A=0\).
\end{lemma}

\begin{proof}
Proposition~\ref{prop:step-006-exact-score-transfer} gives one fixed map
\[
\varphi_A:\mathcal X\to\mathbb R^{r_A}
\]
such that for every \(h\in\mathcal H\), some \(w_h\in\mathbb R^{r_A}\)
satisfies
\[
h(x)\langle w_h,\varphi_A(x)\rangle>0
\qquad\text{for every }x\in\mathcal X.
\]
Thus \(d=r_A\) belongs to the set of admissible dimensions in the definition
of \(\operatorname{dc}(\mathcal H)\), and taking its infimum gives
\(\operatorname{dc}(\mathcal H)\le r_A\).  The admissible set is nonempty,
so its empty-infimum convention is not invoked under these assumptions.

If \(\mathcal H=\varnothing\), dimension zero is admissible: use the unique
map from \(\mathcal X\) to \(\mathbb R^0\), and the target-indexed
requirement is vacuous.  Since admissible dimensions lie in \(\mathbb N_0\),
this gives \(\operatorname{dc}(\mathcal H)=0\).

If \(\mathcal X=\varnothing\), dimension zero is again admissible.  Use the
unique map \(\varnothing\to\mathbb R^0\); for every target, choose the empty
weight, and the pointwise condition is vacuous.  Hence
\(\operatorname{dc}(\mathcal H)=0\).  In this branch,
Proposition~\ref{prop:step-001-mean-response-space} additionally gives
\(r_A=0\).

Finally, suppose \(r_A=0\).  If both \(\mathcal X\) and \(\mathcal H\) were
nonempty, Proposition~\ref{prop:step-003-nonzero-rank} would give
\(r_A\ge1\), a contradiction.  Therefore \(\mathcal X=\varnothing\) or
\(\mathcal H=\varnothing\), and the preceding paragraphs give
\(\operatorname{dc}(\mathcal H)=0=r_A\).  This exhausts the zero-rank
boundary without asking a zero-dimensional score to realize a nonempty sign
constraint.
\end{proof}

\begin{proposition}[Exact conditional polynomial dimension closure]
\label{prop:step-006-polynomial-closure}
Under Assumptions~\ref{assump:parameter-regime}--\ref{assump:mean-response-rank},
and the conclusions of
Proposition~\ref{prop:step-006-exact-score-transfer} and
Lemma~\ref{lem:step-006-dimension-admissibility}, the fixed map
\(\varphi_A:\mathcal X\to\mathbb R^{r_A}\) satisfies, for every
\(h\in\mathcal H\), the existence of \(w_h\in\mathbb R^{r_A}\) with
\[
h(x)\langle w_h,\varphi_A(x)\rangle
\ge1-2\varepsilon=\rho>\frac12>0
\qquad\text{for every }x\in\mathcal X,
\]
and
\[
\operatorname{dc}(\mathcal H)
\le r_A
\le B\bigl(1+m/\tau^2\bigr)^k.
\]
The conclusion is deterministic and contains no hidden constants.  It remains
conditional on the primitive mean-response-rank certificate and makes no
claim about seed-specific terminal rank, a rank-free implication, an
unconditional theorem, or a universal linear bound.
\end{proposition}

\begin{proof}
The pointwise score statement, including pre-instance map independence and
strict positivity, is Proposition~\ref{prop:step-006-exact-score-transfer}.
Lemma~\ref{lem:step-006-dimension-admissibility} gives
\[
\operatorname{dc}(\mathcal H)\le r_A.
\]
Assumption~\ref{assump:mean-response-rank} gives, for the same \(r_A\),
\[
r_A\le B\bigl(1+m/\tau^2\bigr)^k.
\]
Transitivity proves the displayed chain.  Both inequalities are retained
verbatim; no term is absorbed, asymptotic notation introduced, or constant
hidden.

The parameter boundaries follow by direct substitution.  If \(m=0\), then
\[
B\bigl(1+m/\tau^2\bigr)^k=B(1+0)^k=B,
\]
so \(\operatorname{dc}(\mathcal H)\le r_A\le B\).  The unique empty complete
response rule from Lemma~\ref{lem:step-001-complete-rules} applies, and no
query-round argument is used.  If \(\varepsilon=0\), then \(\rho=1\) and
\[
h(x)\langle w_h,\varphi_A(x)\rangle\ge1
\]
with no margin degradation or probabilistic surrogate.  If \(B=1\), the
rank bound is
\[
r_A\le\bigl(1+m/\tau^2\bigr)^k;
\]
if \(k=1\), it is
\[
r_A\le B\bigl(1+m/\tau^2\bigr);
\]
and if \(B=k=1\), it is exactly
\[
r_A\le1+m/\tau^2.
\]
These substitutions use only \(m\in\mathbb N_0\), \(\tau>0\), \(B\ge1\),
and \(k\ge1\), and introduce no upper bound on \(\tau\).

Empty \(\mathcal X\), empty \(\mathcal H\), and \(r_A=0\) are handled by
Lemma~\ref{lem:step-006-dimension-admissibility}.  The primitive certificate,
rather than the number of response transcripts, terminal predictors, or
queries alone, supplies the second dimension inequality.  The probability
mode is deterministic, the horizon mode is the fixed finite upper horizon
\(m\), and the metric consists of exact algebraic dimension and the pointwise
signed score.  Thus every exposed dependence is the one displayed above.
\end{proof}

\subsection{Proof of the Main Theorem}

\begin{proof}[Proof of Theorem~\ref{thm:main}]
Lemma~\ref{lem:step-001-complete-rules} and
Propositions~\ref{prop:step-001-mean-response-space} and
\ref{prop:step-001-fixed-coordinates} construct the fixed all-rule
mean-response space, its finite rank, and one pre-instance coordinate map.
Propositions~\ref{prop:step-002-compact-body} and
\ref{prop:step-002-zero-dimensional}, together with
Lemma~\ref{lem:step-002-continuous-evaluations}, place all subsequent
arguments in one nonempty compact convex body with continuous point
evaluations.

Lemma~\ref{lem:step-003-exact-center-legality} and
Proposition~\ref{prop:step-003-exact-correlation} produce the exact
distributionwise signed-correlation witness, while
Proposition~\ref{prop:step-003-nonzero-rank} settles the nonempty zero-rank
boundary.  Proposition~\ref{prop:step-004-finite-margin} converts those
witnesses, through the fixed-body minimax equality, into an attained
simultaneous margin on every nonempty finite restriction.
Proposition~\ref{prop:step-005-global-margin} then uses the closed constraints
and compact finite-intersection argument to obtain one \(f_h\) on the entire
arbitrary domain.

Finally, Proposition~\ref{prop:step-006-exact-score-transfer} converts
\(f_h\) to exact basis scores with zero residual,
Lemma~\ref{lem:step-006-dimension-admissibility} proves the dimension and
degenerate-case conclusions, and
Proposition~\ref{prop:step-006-polynomial-closure} appends the primitive rank
certificate and the exact parameter specializations.  These statements are
precisely all conclusions of Theorem~\ref{thm:main}.
\end{proof}
